% 从旧数学附录迁出的树路径说明；不新增算法。
\section{树路径交：投影与单点交}\label{knowledge-path-intersection}
\index{树上路径交}
接口见第~\pageref{compact-path_intersection}~页。$x$ 在路径 $a$--$b$ 上的投影是
$\operatorname{lca}(a,b)$、$\operatorname{lca}(a,x)$、$\operatorname{lca}(b,x)$ 中深度最大者（同一根）。
设 $c,d$ 的投影为 $p,q$：$p\ne q$ 时交集为 $p$--$q$；$p=q$ 时若 $d(c,p)+d(p,d)=d(c,d)$ 则为单点交，否则为空。
距离与 depth 按边数，不能换成带权根距离。
返回交集端点和顶点数，空交为 $\{0,0,0\}$，公共边数 $\max(0,\mathrm{vertices}-1)$。
每次调用常数次 LCA；用 HLD 回调为 $O(\log n)$。

\section{非负权树的点集直径摘要}\label{knowledge-set-diameter}
代码见第~\pageref{compact-TreeDiameter}~页。非空点集只需保存一对直径端点 $a,b$，任意点 $x$ 到点集的最远距离为
$\max(d(x,a),d(x,b))$。要求边权非负（可为零），不能推广到一般图或负权树。
插入新点只比较到两个旧端点的距离；合并点集只插入另一摘要的两个端点。并列直径任选，端点不固定。
带权距离可用 $(D_u-D_z)+(D_v-D_z)$（$z$ 为 LCA）；所有路径权和须在 long long 内。

\paragraph{离线添叶仍是固定树度量。}
CF379F：先建出最终树再预处理距离，按出现顺序向摘要插点（初始四点也插入）。$q\le500000$ 时 $4+2q$ 个点，HLD $O(q)$ 空间、$O(q\log q)$ 时间。

QOJ8235（点权互异）：令 $g$ 为全树点权的 mex，按点权 $0,\ldots,g-1$ 建前缀摘要；
某前缀全在以 $x$ 为中心、半径 $k$ 的范围内当且仅当其最远距离 $\le k$，二分最长前缀；$g=0$ 时为零。边带权，不能用边数距离。

\clearpage
